%==================================================
% METODOLOGÍA
%==================================================
\ReportChapter{Metodología}

El presente análisis fue desarrollado mediante un flujo de trabajo estructurado orientado a la preparación, exploración y análisis de datos, así como a la generación de un informe técnico reproducible. La metodología integra herramientas estadísticas, programación y técnicas de visualización para transformar los datos en evidencia útil para la toma de decisiones.



\SectionTitle{Herramientas utilizadas}

Durante el desarrollo del análisis se emplearon las siguientes herramientas:

\begin{itemize}

\item MySQL para la limpieza, transformación y preparación del conjunto de datos.

\item Python como lenguaje principal para el procesamiento y análisis de información.

\item Pandas y NumPy para la manipulación de datos y operaciones estadísticas.

\item Bokeh para la generación de visualizaciones interactivas y gráficos analíticos.

\item Jupyter Notebook como entorno de desarrollo, documentación y ejecución del análisis.

\item \LaTeX{} para la generación automática del informe técnico final.

\end{itemize}



\SectionTitle{Flujo metodológico}

El proceso de análisis se desarrolló mediante las siguientes etapas:


\begin{enumerate}

\item Obtención y revisión inicial del conjunto de datos utilizado en el estudio.

\item Exploración inicial de la información mediante análisis exploratorio de datos (EDA) para evaluar estructura, distribución y características principales de las variables.

\item Preparación y transformación de los datos mediante procesos de limpieza, corrección de formatos y adecuación de variables para el análisis.

\item Validación del conjunto de datos mediante revisiones adicionales para verificar consistencia y disponibilidad de la información procesada.

\item Desarrollo del análisis exploratorio y estadístico orientado a responder las preguntas planteadas en el estudio.

\item Elaboración de tablas, visualizaciones e indicadores utilizados para interpretar los resultados obtenidos.

\item Automatización de la generación de elementos del informe técnico, incluyendo tablas, gráficos y contenido descriptivo.

\item Construcción del informe final mediante \LaTeX{} siguiendo una estructura reproducible y documentada.

\end{enumerate}